% ============================================================================
%  Saussure in HLLSet Algebra --- Part V: Questions of Retrospective Linguistics
%  A computational exegesis of Ferdinand de Saussure's
%  Course in General Linguistics, Part Five, in the Boolean lattice
%  and GF(2) ring of HLLSet Algebra
% ----------------------------------------------------------------------------
%  Series:  Saussure in HLLSet Algebra (reverse direction of influence)
%  Build:   pdflatex saussure_hllset_part5_retrospective_linguistics.tex
%           (run twice)
%  Notation follows Part I:
%    saussure_hllset_part1_general_principles.tex
%  Part II:
%    saussure_hllset_part2_synchronic_linguistics.tex
%  Part III:
%    saussure_hllset_part3_diachronic_linguistics.tex
%  Part IV:
%    saussure_hllset_part4_geographical_linguistics.tex
%  and the EWM PUB series.
% ============================================================================
\documentclass[11pt,a4paper]{article}

\usepackage[utf8]{inputenc}
\usepackage[T1]{fontenc}
\usepackage{amsmath,amssymb,amsthm}
\usepackage{amscd}
\usepackage{booktabs}
\usepackage{array}
\usepackage{hyperref}
\usepackage[margin=1in]{geometry}
\usepackage{enumitem}
\usepackage{xcolor}
\definecolor{hlblue}{RGB}{0,80,160}
\hypersetup{colorlinks=true,linkcolor=hlblue,urlcolor=hlblue,citecolor=hlblue}

% ---- theorems (aligned with the PUB series and Parts I--IV) ----------------
\newtheorem{principle}{Principle}[section]
\newtheorem{theorem}[principle]{Theorem}
\newtheorem{definition}[principle]{Definition}
\newtheorem{proposition}[principle]{Proposition}
\newtheorem{corollary}[principle]{Corollary}
\newtheorem{remark}[principle]{Remark}
\theoremstyle{remark}
\newtheorem{observation}[principle]{Observation}
\newtheorem{reading}[principle]{Reading}

% ---- shorthands (aligned with Parts I--IV and the PUB series) ---------------
\newcommand{\Plane}{\ensuremath{\mathbb{B}}}            % bit plane
\newcommand{\Lat}{\ensuremath{\mathcal{B}}}              % HLLSet (Boolean) lattice
\newcommand{\Tok}{\ensuremath{\mathcal{T}}}              % token space
\newcommand{\HLL}{\ensuremath{H}}                        % an HLLSet
\newcommand{\xor}{\ensuremath{\Delta}}                   % ring addition
\newcommand{\BSS}{\ensuremath{\mathsf{BSS}}}             % BSS measure
\newcommand{\res}{\ensuremath{\mathrm{res}}}             % residual
\newcommand{\spanof}[1]{\ensuremath{\langle #1\rangle}}  % GF(2) span
\newcommand{\supp}{\ensuremath{\mathrm{supp}}}           % support
\newcommand{\Ctx}{\ensuremath{\mathcal{C}}}              % token context
\newcommand{\Ing}{\ensuremath{\mu}}                      % ingestion morphism
\newcommand{\Mat}{\ensuremath{\rho}}                     % materialization morphism
\newcommand{\Lang}{\ensuremath{\mathcal{L}}}             % langue (the lattice system)
\newcommand{\Par}{\ensuremath{\mathcal{P}}}              % parole (usage)
\newcommand{\Sign}{\ensuremath{\mathcal{S}}}             % the sign relation
\newcommand{\Cov}{\ensuremath{\mathrm{cov}}}             % coverage
\newcommand{\Disp}{\ensuremath{\mathcal{D}}}             % the dispatcher
\newcommand{\Sync}{\ensuremath{\mathsf{Syn}}}            % synchronic fiber
\newcommand{\Dia}{\ensuremath{\mathsf{Dia}}}             % diachronic trajectory
\newcommand{\Gen}{\ensuremath{\gamma}}                   % basis generation
\newcommand{\cover}{\ensuremath{\mathrm{cover}}}         % frame cover
\newcommand{\Val}{\ensuremath{\mathcal{V}}}              % value (differential)
\newcommand{\Synt}{\ensuremath{\mathsf{Synt}}}           % syntagmatic reading
\newcommand{\Ass}{\ensuremath{\mathsf{Ass}}}             % associative reading
\newcommand{\Dep}{\ensuremath{\mathsf{Dep}}}             % departed
\newcommand{\Ret}{\ensuremath{\mathsf{Ret}}}             % retained
\newcommand{\New}{\ensuremath{\mathsf{New}}}             % new
\newcommand{\Trans}{\ensuremath{\mathsf{Tr}}}            % recorded transition
\newcommand{\Piv}{\ensuremath{\mathsf{Piv}}}             % pivot
\newcommand{\Ana}{\ensuremath{\mathsf{Ana}}}             % analogical fourth
\newcommand{\Change}{\ensuremath{\mathbf{M}}}            % change-of-basis matrix
\newcommand{\World}{\ensuremath{\mathsf{W}}}             % a world
\newcommand{\Iota}{\ensuremath{\iota}}                   % interpretation map
\newcommand{\Bridge}{\ensuremath{\Phi}}                  % bridge between worlds
\newcommand{\Wset}{\ensuremath{\Omega}}                  % family of worlds
% ---- new shorthands for Part V ---------------------------------------------
\newcommand{\Recon}{\ensuremath{\mathsf{R}}}             % reconstruction operator
\newcommand{\Proto}{\ensuremath{\mathsf{P}}}             % proto-node (reconstruction)
\newcommand{\Fam}{\ensuremath{\mathsf{Fam}}}             % language family
\newcommand{\Typ}{\ensuremath{\mathsf{Typ}}}             % linguistic type

% ---- title -----------------------------------------------------------------
\title{\textbf{Saussure in HLLSet Algebra}\\[4pt]
\large Part V: Questions of Retrospective Linguistics\\[4pt]
\large A Computational Exegesis of the \textit{Course in General Linguistics},\\
Part Five, in the Boolean Lattice and GF(2) Ring of HLLSet Algebra\\[6pt]
\normalsize Fifth paper in the series \textit{Saussure in HLLSet Algebra}}

\author{Alex Mylnikov \and Alexandr Mylnikov Jr. \and Deependra Kumar}
\date{\today}

\begin{document}
\maketitle

\begin{abstract}
Part V closes the series by re-founding Part Five of the \textit{Course},
\emph{Questions of Retrospective Linguistics}, as time travel through old
frames. The two perspectives of diachronic linguistics become the two
directions of the commit DAG: the prospective reading follows the pivots
forward, the retrospective reading inverts the transport matrices backward.
The inversion is exact on the old span---the map $c \mapsto c\Change$ is
injective---and partial elsewhere, where the spill against the earliest frame
measures the error of reading the present through the past's lens. The search
for the earliest language and the prototype is re-founded as a fact about
perishable origins: the lineage is not monotone, its frame transitions
depart, retain, and add bits, and the first basis is an arbitrary origin
that can be entirely replaced; the prototype is therefore a backward path
along the lineage, not a projection into the first basis. Reconstruction
becomes an executable
inverse problem: the proto-node of a set of descendants is their meet,
$\Proto = \bigcap_i X_i$, a view whose provenance points into the past; the
reconstruction is exact exactly when nothing has departed, and the spill of
$\Proto$ against the earliest frame measures what the descendants can no
longer witness. The contribution of linguistic evidence to anthropology and
prehistory is delimited with the same discipline: the DAG certifies shared
ancestry of frames and shared nodes, and nothing else. Language families and
linguistic types are then separated: a family is ancestry in the DAG, a type
is similarity in the current fiber, and the algebra computes both without
forcing them to agree. The paper, and the series, close with Saussure's own
conclusion re-founded: the unique and true object of linguistics is langue
envisaged in itself and for itself---and the lattice is exactly that object.
\end{abstract}

\noindent\textbf{Keywords:} Saussure; retrospective linguistics; prospective
method; reconstruction; prototype; language family; linguistic type; time
travel; old frames; transport inversion; meet; provenance; HLLSet; Boolean
lattice; GF(2) ring; computational exegesis.

% ============================================================================
\section{Introduction}
\label{sec:intro}
% ============================================================================

\subsection{Where Part V stands in the series}

Part I \cite{part1} re-founded the general principles; Part II \cite{part2}
the synchronic theory of value; Part III \cite{part3} the diachronic theory
of the pivot; Part IV \cite{part4} the geographical theory of worlds and
waves. Part Five of the \textit{Course}, Questions of Retrospective
Linguistics, is where Saussure turns the method around: instead of following
the course of time, the linguist goes backward, reconstructing what the
present can no longer show directly \cite[Part V]{harris1983}. Part V of this
series re-founds that backward method in the algebra, and with it closes the
series.

\subsection{Method, unchanged}

The method is that of Parts I--IV \cite{part1,part2,part3,part4}: no
metaphor, principles as theorems or invariants, ambiguity recorded as data.
The translation note of Part I applies unchanged; references to the
\textit{Course} are by part, chapter, and section.

\subsection{Roadmap}

Section~\ref{sec:prelim} fixes the standing definitions and adds the
retrospective layer (prospective and retrospective readings, the
reconstruction operator). Section~\ref{sec:ch1} re-founds the two
perspectives of diachronic linguistics as the two directions of the DAG.
Section~\ref{sec:ch2} re-founds the search for the earliest language and the
prototype as a backward path along a perishable lineage, not a projection
into the first basis. Section~\ref{sec:ch3} is the
center: reconstruction as the meet of descendants, with its exactness
criterion. Section~\ref{sec:ch4} delimits the anthropological evidence.
Section~\ref{sec:ch5} separates families from types. Section
\ref{sec:conclusion} closes the series with Saussure's own conclusion
re-founded. Section~\ref{sec:dictionary} collects the dictionary of Part
Five.

% ============================================================================
\section{Preliminaries: the whole inheritance and the retrospective layer}
\label{sec:prelim}
% ============================================================================

We restate the standing definitions with their labels; proofs and motivation
are in \cite{part1,part2,part3,part4}.

\begin{definition}[Standing definitions]
\label{def:standing}
The \textbf{plane} is $\Plane=\{0,\dots,32767\}$; the \textbf{lattice} is
$\Lat=\mathcal{P}(\Plane)$ under $\cup,\cap,\subseteq$ and the Boolean ring
$\xor=\Delta$, $\cap$. $\Ing:\Tok\to\Lat$ is an \textbf{IICA morphism};
$\Mat:\Lat\dashrightarrow\Ctx$ materializes through the door. A
\textbf{basis frame} $F=(B_1,\dots,B_k)$ is a row-echelon GF(2) basis with
distinct leading bits, stamped with a generation $\Gen$; hard keys, soft
keys, and spill are as in Parts I--III. The \textbf{synchronic fiber} is
$\Sync(t)=(\Lat,F(t),\Gen(t))$; a \textbf{pivot} is the extension/rotation
of Part III; the \textbf{transport} across a pivot is $c' = c\Change$ with
the change-of-basis matrix $\Change$ \cite{part3}. A \textbf{world} and a
\textbf{bridge} are as in Part IV \cite{part4}.
\end{definition}

\subsection{The retrospective layer}

\begin{definition}[Prospective and retrospective readings]
\label{def:perspectives}
The \textbf{prospective reading} of a lineage is the forward traversal of the
commit DAG: at each pivot, transport coordinates forward,
$c(t{+}1) = c(t)\Change(t{+}1)$. The \textbf{retrospective reading} is the
backward traversal: given the present coordinates $c(t)$, recover the past
coordinates $c(t{-}1)$ of any node that was already in the old span, and
project nodes outside the old span onto the old frame with a measured error.
\end{definition}

\begin{definition}[Reconstruction operator]
\label{def:recon}
Let $X_1,\dots,X_n$ be nodes (descendants) in the current lattice. The
\textbf{reconstruction} of their common ancestor is the meet
\[
  \Recon(X_1,\dots,X_n)
  \;=\;
  \Proto
  \;=\;
  \bigcap_{i=1}^n X_i,
\]
read as a \emph{view}: a recipe whose sources are the descendants and whose
provenance points into the past. The \textbf{reconstruction spill} is
\[
  s(\Proto) \;=\; |\Proto \setminus \cover(F_0 \cap F_t)|
\]
against the \emph{surviving} part of the earliest recorded frame $F_0$;
the departed part $F_0 \setminus F_t$ is accounted as uncertainty
(Definition~\ref{def:frame-transition}), not as a lens.
\end{definition}

\begin{proposition}[Backward reading is exact on the old span]
\label{prop:backward}
Let $F(t{-}1)=(B_1,\dots,B_k)$ and $F(t)=(B'_1,\dots,B'_{k'})$ be two frames
of a monotone ring, $k' \ge k$, with change-of-basis matrix $\Change$
($k \times k'$). Then:
\begin{enumerate}[leftmargin=*,nosep]
  \item the rows of $\Change$ are the coordinates of the old basis in the new
    basis; they are linearly independent, so the map
    $c \mapsto c\Change$ is injective;
  \item consequently, for every node $X \in \spanof{F(t{-}1)}$, the past
    coordinates are uniquely recoverable from the present coordinates: the
    equation $c' = c\Change$ has at most one solution $c$, and it has one
    exactly when $c'$ lies in the image (i.e.\ when $X$ was in the old
    span);
  \item for $X \notin \spanof{F(t{-}1)}$, no past coordinates exist, and the
    old-frame projection has spill
    $|X \setminus \cover(F(t{-}1))|$, the measured error of reading the
    present through the past's lens.
\end{enumerate}
\end{proposition}

\begin{proof}
(1) The old basis vectors are linearly independent over GF(2), so their
coordinate vectors are linearly independent; hence $\Change$ has full row
rank $k$, and left multiplication by $\Change$ is injective. (2) Immediate
from (1). (3) A node outside the old span has no old hard key; the spill is
the part of it that the old frame cannot see, as in Part III
\cite{part3,state2026}.
\end{proof}

\begin{remark}[Monotone span versus perishable frames]
\label{rem:monotone}
Proposition~\ref{prop:backward} states the retrospective method in the
monotone-span regime of the EWM ring, where
$\spanof{F(t{-}1)} \subseteq \spanof{F(t)}$ and no direction is ever lost.
A real lineage is not monotone at the frame level: each new frame differs
from its predecessor by a frame transition with departed, retained, and new
parts, and the departed bits accumulate---at the limit they can replace the
first frame entirely. In that regime the backward map is exact only on the
retained part of the old span; the departed part is not inverted but
reconstructed by hypothesis. This is the subject of Chapters II and III.
\end{remark}

% ============================================================================
\section{Chapter I: The two perspectives of diachronic linguistics}
\label{sec:ch1}
% ============================================================================

Saussure distinguishes two ways of doing diachronic linguistics: the
\emph{prospective} method follows the course of time from past documents to
the present; the \emph{retrospective} method goes backward from the present
to the past, and it is the method of reconstruction \cite[Part V,
ch.~I]{harris1983}.

\begin{reading}[The two directions of the DAG]
\label{read:perspectives}
The prospective perspective is the forward traversal of the commit DAG: it
sees the pivots as they happen, and transports coordinates forward. The
retrospective perspective is the backward traversal: it sees the present
fiber and asks which earlier fiber could have produced it, inverting
transports by Proposition~\ref{prop:backward}.
\end{reading}

\begin{observation}[The two methods have different exactness]
The prospective reading is exact by construction: every pivot and every
transport is recorded. The retrospective reading is exact only on the old
span, and partial outside it, where the spill measures the error. Saussure's
distinction between the two methods is therefore, in the algebra, a
distinction between forward determinism and backward inversion with a
measured residual---not a symmetrical pair of viewpoints.
\end{observation}

% ============================================================================
\section{Chapter II: Earliest languages and prototypes}
\label{sec:ch2}
% ============================================================================

Saussure's chapter II deflates the search for the first language: no living
language is ``older'' than another in any absolute sense, and the prototype
reconstructed by comparison is a hypothesis about relations, not a reality
that was ever spoken \cite[Part V, ch.~II]{harris1983}.

\begin{definition}[Frame transitions along a lineage]
\label{def:frame-transition}
Let $F_0, F_1, \dots, F_t$ be the recorded frames of a lineage. Each step is
a \textbf{frame transition} with departed, retained, and new parts
\[
  \Dep_i = F_{i-1}\setminus F_i, \qquad
  \Ret_i = F_{i-1}\cap F_i, \qquad
  \New_i = F_i\setminus F_{i-1},
\]
which partition $F_{i-1}\cup F_i$. The \textbf{accumulated departure} from
$F_0$ to $F_t$ is
\[
  \Dep(0,t) \;=\; F_0 \setminus F_t .
\]
Nothing prevents $\Dep(0,t)$ from growing until it equals $F_0$: the
lineage is not monotone, and the first frame can be entirely replaced.
\end{definition}

\begin{reading}[The first frame is an arbitrary and perishable origin]
\label{read:origin}
The earliest frame $F_0$ of a lineage is the first committed generation in
the DAG. It is an \emph{origin} only because the record starts there, not
because the plane has an origin: by the coordinate-freedom theorem of Part I
\cite{part1}, any re-coordination of the plane is an automorphism, and no
structural property prefers one first basis over another. And the origin is
\emph{perishable}: the lineage proceeds by frame transitions
(Definition~\ref{def:frame-transition}), whose departed parts accumulate.
The prototype is therefore not a projection of the present into $F_0$---a
lens that may have vanished---but a backward path along the lineage.
\end{reading}

\begin{observation}[No language is older than another]
\label{obs:no-oldest}
For every lineage, the first commit is arbitrary: it is the point where the
record begins, and a different beginning would be an isomorphic coordinate
choice. The question ``which language is oldest'' therefore has no
algebraic meaning: all lineages are equally old, because age is a property
of the record, not of the lattice. Saussure's deflation of the question is
re-founded as the arbitrariness of the origin---an origin that, in
addition, does not survive its own lineage.
\end{observation}

\begin{proposition}[The prototype is a backward path, not a projection]
\label{prop:prototype}
Let $F_0,\dots,F_t$ be the frames of a lineage and $X$ a present node.
\begin{enumerate}[leftmargin=*,nosep]
  \item A single projection of $X$ into $F_0$ is not the prototype: it
    reads the present through an origin that may have departed in part or
    in whole. When $\Dep(0,t) = F_0$, the lens $F_0$ is foreign to the
    present frame ($F_0 \cap F_t = \varnothing$), and the projection
    measures similarity to a vanished system, not descent from it.
  \item The correct prototype reading of $X$ is the backward path along
    the lineage: transport $X$'s coordinates backward through each
    \emph{retained} part $\Ret_i$, and supply the \emph{departed} parts
    $\Dep_i$ only by the comparative hypothesis of Section~\ref{sec:ch3},
    applied now to oppositions rather than to nodes.
  \item The accumulated departure $\Dep(0,t)$ is the exact uncertainty of
    the prototype: the part of the origin that the present can no longer
    witness, and that only hypothesis---never the record---can supply.
\end{enumerate}
\end{proposition}

% ============================================================================
\section{Chapter III: Reconstructions}
\label{sec:ch3}
% ============================================================================

Saussure's chapter III is the heart of Part Five: reconstruction is
comparative, the starred form is a hypothesis about \emph{relations}, and its
value lies in the correspondences it summarizes, not in any claim about
substance \cite[Part V, ch.~III]{harris1983}.

\begin{reading}[Reconstruction is a meet]
\label{read:recon}
The comparative method places descendant forms side by side and asks what
they have in common. In the lattice, the descendants are nodes
$X_1,\dots,X_n$, and what they have in common is the largest node contained
in all of them---the meet $\Proto = \bigcap_i X_i$. The starred form is
$\Proto$ read as a view: a recipe over the descendants' provenance, stamped
with the present generation, and explicitly \emph{not} an original.
\end{reading}

\begin{proposition}[Exactness of reconstruction]
\label{prop:recon}
Suppose the true ancestor $A$ exists in the DAG and each descendant
$X_i$ was formed from $A$ by joins (additions of bits), so $A \subseteq X_i$
for all $i$. Then:
\begin{enumerate}[leftmargin=*,nosep]
  \item $A \subseteq \Proto$: the reconstruction never claims less than the
    ancestor;
  \item $\Proto \setminus A$ is the innovation shared by \emph{all}
    descendants (the convergent addition), which reconstruction cannot
    distinguish from inheritance;
  \item $A \setminus \Proto$ is the part of the ancestor that departed in at
    least one descendant, which the descendants alone cannot recover;
  \item reconstruction is exact ($\Proto = A$) if and only if nothing
    departed, i.e.\ $A \setminus \Proto = \varnothing$.
\end{enumerate}
The reconstruction spill of Definition~\ref{def:recon} measures the part of
$\Proto$ that the \emph{surviving} part of the earliest frame cannot see;
the departed part of the earliest frame is itself uncertainty. The
differences $\Proto \setminus A$ and $A \setminus \Proto$ are the two
directions of the reconstruction's uncertainty.
\end{proposition}

\begin{proof}
(1) $A \subseteq X_i$ for all $i$ implies $A \subseteq \bigcap_i X_i =
\Proto$. (2) If a bit lies in every descendant but not in $A$, it was added
convergently in all lines; the meet cannot tell addition from inheritance.
(3) A bit of $A$ missing from some descendant is missing from the meet. (4)
Immediate from (1)--(3).
\end{proof}

\begin{corollary}[Correspondences are transported coordinate agreements]
\label{cor:correspondences}
When the descendants live in different worlds or different generations, the
comparative method first transports their hard keys to a common reference
frame (an ancestor frame if one is shared); a \emph{correspondence} between
two descendants is the equality of transported coordinates on a set of
oppositions. The reconstruction $\Proto$ is then the node whose transported
coordinates are the meet of the descendants' coordinates---the exact summary
of the correspondences, and nothing more.
\end{corollary}

\begin{remark}[The starred form is a hypothesis, not a sound]
Saussure insists that a reconstruction is a set of relations, not a claim
about how the proto-language sounded. The algebra agrees precisely: $\Proto$
is a view---an expression over the descendants---and it is never stored as an
original. Its generation stamp is the present generation in which the
hypothesis is evaluated, and its provenance points into the past without
claiming that the node itself was ever an observed original.
\end{remark}

% ============================================================================
\section{Chapter IV: Linguistic evidence in anthropology and prehistory}
\label{sec:ch4}
% ============================================================================

Saussure's chapter IV asks what linguistics can contribute to anthropology
and prehistory, and answers with characteristic discipline: linguistic
evidence is structural, and it must not be made to prove what only other
sciences can prove \cite[Part V, ch.~IV]{harris1983}.

\begin{reading}[The DAG certifies ancestry, nothing else]
\label{read:anthropology}
The algebra's evidence is of exactly three kinds: shared nodes (cognates),
shared frames (a common system of oppositions), and shared pivots (a common
history of changes). From these it can certify common ancestry of frames and
the order of divergences. It cannot certify anything about the speakers'
genes, pots, or migrations, because bits, generations, and provenance carry
no such information.
\end{reading}

\begin{observation}[The discipline of the lattice]
\label{obs:discipline}
The lattice disciplines the inference in both directions. It \emph{forces}
the linguist to say exactly what is shared---a node, a frame, a pivot---and
it \emph{forbids} conclusions that the shared objects do not support.
Saussure's caution about linguistic evidence is therefore not a stylistic
preference; it is a type discipline: an inference is valid only if it is a
function of the DAG, the worlds, and the provenance, and nothing else.
\end{observation}

% ============================================================================
\section{Chapter V: Language families and linguistic types}
\label{sec:ch5}
% ============================================================================

Saussure's last chapter separates two classifications that are constantly
confused: the \emph{family} is genealogical, the \emph{type} is structural.
They can and do diverge \cite[Part V, ch.~V]{harris1983}.

\begin{reading}[Families are ancestry; types are profiles]
\label{read:family-type}
A family is a property of the DAG: two worlds belong to the same family when
their frames descend from a common ancestor frame. A type is a property of
the current fiber: two worlds are of the same type when their soft-key
profiles against a reference frame are close. The two computations use
disjoint data and need not agree.
\end{reading}

\begin{definition}[Family and type]
\label{def:family-type}
The \textbf{family} of a world $\World$ is
\[
  \Fam(\World) \;=\;\{\,\World' : F(\World) \text{ and } F(\World')
  \text{ descend from a common ancestor frame}\,\},
\]
an ancestry relation in the frame DAG. The \textbf{type} of a world against
a reference frame $F$ is its profile
\[
  \Typ_F(\World) \;=\;\bigl(\BSS_{B_1}(\mathcal{N}(\World)),\dots,
  \BSS_{B_k}(\mathcal{N}(\World))\bigr),
\]
a similarity relation in the current fiber. Family is a fact of provenance;
type is a fact of the present.
\end{definition}

\begin{proposition}[Families and types are independent]
\label{prop:family-type}
Two unrelated worlds may be type-identical (convergence): their histories
differ but their present profiles coincide. Two related worlds may be
type-distant (divergence): their histories coincide up to a point but their
profiles have drifted apart. The algebra computes both and never forces them
to agree; Saussure's warning against confusing family and type is re-founded
as the independence of the ancestry relation and the profile metric.
\end{proposition}

\begin{proof}
Ancestry is determined by the DAG; type is determined by BSS against a
fixed frame. The two functions have different domains (history vs present)
and different data (provenance vs soft keys); nothing in the definitions
couples them.
\end{proof}

% ============================================================================
\section{Conclusion of the series: langue in itself and for itself}
\label{sec:conclusion}
% ============================================================================

Saussure closes the \textit{Course} with the sentence that defined a
century: linguistics has for its unique and true object \emph{la langue
envisag\'ee en elle-m\^eme et pour elle-m\^eme}---language studied in itself
and for itself \cite[Conclusion]{harris1983}. The series closes by
re-founding that sentence, and it does so with an exactness that no
philosophical commentary could offer.

\emph{In itself.} The lattice $\Lat$ over the plane $\Plane$ is studied in
itself: its Boolean structure and its GF(2) ring, its frames, generations,
pivots, and transports. Nothing outside the lattice is needed to state its
laws; the panchronic axioms of Part I are the exact form of ``in itself.''

\emph{For itself.} The lattice is not the servant of any host, any
vocabulary, any speaker. It is the system of differences that Parts I--V
re-founded chapter by chapter: the sign as $(\Ing,\Mat)$, arbitrariness as
coordinate freedom, value as the hard key, the pivot as the diachronic
event, analogy as the fourth proportional, worlds as languages, bridges as
waves, reconstruction as the meet. The algebra interprets the
\textit{Course}; the \textit{Course} interprets the algebra; and the object
of both is one and the same: langue, as a lattice, in itself and for itself.

\begin{table}[htbp]
\centering
\renewcommand{\arraystretch}{1.15}
\begin{tabular}{@{}>{\raggedright\arraybackslash}p{3.2cm}
                    >{\raggedright\arraybackslash}p{4.2cm}
                    >{\raggedright\arraybackslash}p{7.2cm}@{}}
\toprule
Part of the \textit{Course} & In the \textit{Course} & In HLLSet Algebra (this series) \\
\midrule
Introduction & langue and parole; the object of linguistics & the morphism
step and the lattice as langue \cite{langue2026}; folded into Part II \\
Part One & general principles & sign, arbitrariness, linearity,
immutability/mutability, synchrony/diachrony \cite{part1} \\
Part Two & synchronic linguistics & value as the hard key; differences
without positive terms; syntagmatic/associative readings \cite{part2} \\
Part Three & diachronic linguistics & the pivot; sound change, analogy, folk
etymology, agglutination; diachronic units \cite{part3} \\
Part Four & geographical linguistics & worlds, codebooks, thresholds;
diversity, continua, substrata, waves \cite{part4} \\
Part Five & retrospective linguistics; conclusion & time travel through old
frames; reconstruction as the meet; families vs types; langue in itself and
for itself \cite{part5} \\
\bottomrule
\end{tabular}
\caption{The complete architecture of the series: the \textit{Course} in
General Linguistics, re-founded part by part in HLLSet Algebra.}
\label{tab:series}
\end{table}

% ============================================================================
\section{The dictionary of Part Five}
\label{sec:dictionary}
% ============================================================================

\begin{table}[htbp]
\centering
\renewcommand{\arraystretch}{1.2}
\begin{tabular}{@{}>{\raggedright\arraybackslash}p{3.4cm}
                    >{\raggedright\arraybackslash}p{4.2cm}
                    >{\raggedright\arraybackslash}p{7.0cm}@{}}
\toprule
Saussure (Part Five) & In the \textit{Course} & In HLLSet Algebra \\
\midrule
prospective method & follow the course of time & forward traversal of the
DAG; transports $c' = c\Change$: Definition~\ref{def:perspectives} \\
retrospective method & go backward from the present & backward inversion on
the old span, projection with spill elsewhere:
Proposition~\ref{prop:backward} \\
earliest language & no language is older than another & the first commit is
an arbitrary origin: Observation~\ref{obs:no-oldest} \\
prototype & a hypothesis, not a reality & a backward path along the lineage;
the first basis is arbitrary and perishable, and $\Dep(0,t)$ is the
uncertainty: Proposition~\ref{prop:prototype} \\
reconstruction & comparative; the starred form & the meet
$\Proto = \bigcap_i X_i$, a view with past provenance:
Definition~\ref{def:recon}, Proposition~\ref{prop:recon} \\
correspondences & the relations reconstruction summarizes & transported
coordinate agreements: Corollary~\ref{cor:correspondences} \\
linguistic evidence in prehistory & structural, disciplined & shared nodes,
frames, pivots certify ancestry, nothing else:
Observation~\ref{obs:discipline} \\
language family & genealogical & ancestry in the frame DAG:
Definition~\ref{def:family-type} \\
linguistic type & structural similarity & soft-key profile in the current
fiber: Definition~\ref{def:family-type} \\
family vs type & they can diverge & independence of ancestry and profile:
Proposition~\ref{prop:family-type} \\
the conclusion of the \textit{Course} & langue in itself and for itself & the
lattice in itself and for itself: Section~\ref{sec:conclusion} \\
\bottomrule
\end{tabular}
\caption{The complete dictionary of Part Five of the \textit{Course} into
HLLSet Algebra.}
\label{tab:dictionary}
\end{table}

% ============================================================================
\section{What Part V adds, and what the series as a whole adds}
\label{sec:adds}
% ============================================================================

\begin{enumerate}[leftmargin=*,nosep]
  \item \textbf{Backward reading becomes a computation.}
    Proposition~\ref{prop:backward} gives the retrospective method its exact
    domain: injective inversion on the old span, measured spill outside it.
  \item \textbf{Reconstruction becomes an executable inverse problem.}
    Definition~\ref{def:recon} and Proposition~\ref{prop:recon} turn the
    starred form into a meet with a proven exactness criterion and a measured
    uncertainty in both directions ($\Proto \setminus A$ and
    $A \setminus \Proto$).
  \item \textbf{The prototype is demystified.} The first basis is an
    arbitrary origin by Part I's coordinate freedom, and it is perishable:
    the lineage proceeds by frame transitions whose departed parts
    accumulate. The prototype is therefore a backward path along the
    lineage, with the accumulated departure as its uncertainty---not a
    projection into a vanished origin.
  \item \textbf{Family and type are separated as datatypes.}
    Definition~\ref{def:family-type} and Proposition~\ref{prop:family-type}
    make the genealogical/structural distinction a difference of data, not
    of emphasis.
  \item \textbf{The series closes the circle.} The \textit{Course} has been
    read out of the algebra in five parts; the algebra has been read out of
    the \textit{Course} in the five EWM papers \cite{ewm2026,companion2026,
    state2026,attention2026,langue2026}; the two directions meet in one
    construction, and the closing object is Saussure's own: langue, in
    itself and for itself.
\end{enumerate}

\section*{Acknowledgments}
This paper was created in collaboration with, and with the assistance of, an
AI system---DeepSeek (deepseek-v4-pro). The meet reading of reconstruction,
the injectivity reading of the retrospective method, the family/type
separation, and several of the formulations emerged in dialogue with the
assistant, whose drafts the authors revised and take full responsibility for
in their final form.

\begin{thebibliography}{99}

\bibitem{saussure1916} F.~de Saussure,
\textit{Cours de linguistique g\'en\'erale}, Payot, Paris, 1916.
(Edited posthumously by C.~Bally and A.~Sechehaye, with A.~Riedlinger.)

\bibitem{harris1983} F.~de Saussure,
\textit{Course in General Linguistics}, translated and annotated by
Roy Harris, Duckworth, London, 1983.

\bibitem{baskin1959} F.~de Saussure,
\textit{Course in General Linguistics}, translated by Wade Baskin,
Philosophical Library, New York, 1959.

\bibitem{saussure1977} F.~de Sossyur,
\textit{Trudy po yazykoznaniyu} [Works in Linguistics], edited by
A.~A.~Kholodovich, translation of the \textit{Cours} by A.~M.~Sukhotin,
Progress, Moscow, 1977.
(Authors' working copy; used for verification of the readings.)

\bibitem{part1} A.~Mylnikov, A.~Mylnikov Jr., and D.~Kumar,
\textit{Saussure in HLLSet Algebra, Part I: General Principles},
SAUSSURE PUB series, 2026.

\bibitem{part2} A.~Mylnikov, A.~Mylnikov Jr., and D.~Kumar,
\textit{Saussure in HLLSet Algebra, Part II: Synchronic Linguistics},
SAUSSURE PUB series, 2026.

\bibitem{part3} A.~Mylnikov, A.~Mylnikov Jr., and D.~Kumar,
\textit{Saussure in HLLSet Algebra, Part III: Diachronic Linguistics},
SAUSSURE PUB series, 2026.

\bibitem{part4} A.~Mylnikov, A.~Mylnikov Jr., and D.~Kumar,
\textit{Saussure in HLLSet Algebra, Part IV: Geographical Linguistics},
SAUSSURE PUB series, 2026.

\bibitem{part5} A.~Mylnikov, A.~Mylnikov Jr., and D.~Kumar,
\textit{Saussure in HLLSet Algebra, Part V: Questions of Retrospective
Linguistics}, SAUSSURE PUB series, 2026.

\bibitem{ewm2026} A.~Mylnikov, A.~Mylnikov Jr., and D.~Kumar,
\textit{Emerging World Models: A Structural Framework for Self-Organizing
World Representations}, EWM PUB series, 2026.

\bibitem{companion2026} A.~Mylnikov, A.~Mylnikov Jr., and D.~Kumar,
\textit{The EWM Lattice and Its Worlds}, EWM PUB series, 2026.

\bibitem{state2026} A.~Mylnikov, A.~Mylnikov Jr., and D.~Kumar,
\textit{Transformations and Trajectories of a Content-Addressed State
Machine}, EWM PUB series, 2026.

\bibitem{attention2026} A.~Mylnikov, A.~Mylnikov Jr., and D.~Kumar,
\textit{Attention on Lattice Traversal Is All You Need}, EWM PUB series, 2026.

\bibitem{langue2026} A.~Mylnikov, A.~Mylnikov Jr., and D.~Kumar,
\textit{Langue as a Lattice: Closing the Circle from Statistical to
Structural Linguistics through Content-Addressed Algebras}, EWM PUB series,
2026.

\bibitem{birkhoff1940} G.~Birkhoff,
\textit{Lattice Theory}, American Mathematical Society, 1940.

\bibitem{maclane1971} S.~Mac Lane,
\textit{Categories for the Working Mathematician}, Springer, 1971.

\bibitem{hjelmslev1961} L.~Hjelmslev,
\textit{Prolegomena to a Theory of Language}, University of Wisconsin Press,
1961.

\bibitem{harris1954} Z.~S.~Harris,
\textit{Distributional Structure}, Word, 10(2--3):146--162, 1954.

\bibitem{firth1957} J.~R.~Firth,
\textit{A Synopsis of Linguistic Theory, 1930--1955}, Studies in Linguistic
Analysis, 1957.

\bibitem{vaswani2017} A.~Vaswani, N.~Shazeer, N.~Parmar, J.~Uszkoreit,
L.~Jones, A.~N.~Gomez, \L.~Kaiser, and I.~Polosukhin,
\textit{Attention Is All You Need}, NeurIPS, 2017.

\end{thebibliography}

\end{document}
